\documentclass[11pt]{article}
\newcommand{\videoid}{V3OboR_tSEM}
% Preámbulo compartido de las notas del curso de topología algebraica.
% Cada nota empieza con \documentclass[11pt]{article}, define \videoid{...}
% (el identificador del video de YouTube) y luego hace % Preámbulo compartido de las notas del curso de topología algebraica.
% Cada nota empieza con \documentclass[11pt]{article}, define \videoid{...}
% (el identificador del video de YouTube) y luego hace % Preámbulo compartido de las notas del curso de topología algebraica.
% Cada nota empieza con \documentclass[11pt]{article}, define \videoid{...}
% (el identificador del video de YouTube) y luego hace \input{preambulo}.
% Compilar con LuaLaTeX:  lualatex nombre-de-la-nota.tex
\usepackage[spanish,provide=*]{babel}
\usepackage[T1]{fontenc}
\usepackage{lmodern}
\usepackage[a4paper,margin=2.6cm]{geometry}
\usepackage{amsmath,amssymb,amsthm}
\usepackage{xcolor}
\usepackage{tikz}
\usetikzlibrary{arrows.meta,babel}
\usepackage[colorlinks=true,urlcolor=blue!50!black,linkcolor=blue!50!black]{hyperref}

\hyphenpenalty=10000 \exhyphenpenalty=10000 \emergencystretch=3em

% Colores de las figuras (los mismos que usa la página)
\definecolor{extremo}{RGB}{47,107,79}
\definecolor{azul}{RGB}{44,74,154}

% Definiciones numeradas por unidad
\theoremstyle{definition}
\newtheorem{defn}{Definición}
% Numeración por unidad: \unidad{u}{n} (después de \input{preambulo}) hace que las
% definiciones de la nota se llamen u.n, u.(n+1), ...; n es el primer número
% libre de esa unidad (ver la tabla de CONVENCIONES.md).
\newcommand{\unidad}[2]{\renewcommand{\thedefn}{#1.\arabic{defn}}\setcounter{defn}{\numexpr#2-1\relax}}

% Observaciones y erratas del video. Usar SOLO si hay algo que señalar:
%   \begin{erratas}
%     \errata{203}{3:23}{Se dice ..., pero debe decir ...}
%   \end{erratas}
% El primer argumento son los segundos (enlaza al minuto exacto del video),
% el segundo es la marca visible, el tercero es el texto.
\newenvironment{erratas}{\section*{Observaciones y erratas del video}\begin{itemize}}{\end{itemize}}
\newcommand{\errata}[3]{%
  \edef\errurl{https://www.youtube.com/watch?v=\videoid\string&t=#1}%
  \item[\textbf{\expandafter\href\expandafter{\errurl}{#2}}] #3}
.
% Compilar con LuaLaTeX:  lualatex nombre-de-la-nota.tex
\usepackage[spanish,provide=*]{babel}
\usepackage[T1]{fontenc}
\usepackage{lmodern}
\usepackage[a4paper,margin=2.6cm]{geometry}
\usepackage{amsmath,amssymb,amsthm}
\usepackage{xcolor}
\usepackage{tikz}
\usetikzlibrary{arrows.meta,babel}
\usepackage[colorlinks=true,urlcolor=blue!50!black,linkcolor=blue!50!black]{hyperref}

\hyphenpenalty=10000 \exhyphenpenalty=10000 \emergencystretch=3em

% Colores de las figuras (los mismos que usa la página)
\definecolor{extremo}{RGB}{47,107,79}
\definecolor{azul}{RGB}{44,74,154}

% Definiciones numeradas por unidad
\theoremstyle{definition}
\newtheorem{defn}{Definición}
% Numeración por unidad: \unidad{u}{n} (después de % Preámbulo compartido de las notas del curso de topología algebraica.
% Cada nota empieza con \documentclass[11pt]{article}, define \videoid{...}
% (el identificador del video de YouTube) y luego hace \input{preambulo}.
% Compilar con LuaLaTeX:  lualatex nombre-de-la-nota.tex
\usepackage[spanish,provide=*]{babel}
\usepackage[T1]{fontenc}
\usepackage{lmodern}
\usepackage[a4paper,margin=2.6cm]{geometry}
\usepackage{amsmath,amssymb,amsthm}
\usepackage{xcolor}
\usepackage{tikz}
\usetikzlibrary{arrows.meta,babel}
\usepackage[colorlinks=true,urlcolor=blue!50!black,linkcolor=blue!50!black]{hyperref}

\hyphenpenalty=10000 \exhyphenpenalty=10000 \emergencystretch=3em

% Colores de las figuras (los mismos que usa la página)
\definecolor{extremo}{RGB}{47,107,79}
\definecolor{azul}{RGB}{44,74,154}

% Definiciones numeradas por unidad
\theoremstyle{definition}
\newtheorem{defn}{Definición}
% Numeración por unidad: \unidad{u}{n} (después de \input{preambulo}) hace que las
% definiciones de la nota se llamen u.n, u.(n+1), ...; n es el primer número
% libre de esa unidad (ver la tabla de CONVENCIONES.md).
\newcommand{\unidad}[2]{\renewcommand{\thedefn}{#1.\arabic{defn}}\setcounter{defn}{\numexpr#2-1\relax}}

% Observaciones y erratas del video. Usar SOLO si hay algo que señalar:
%   \begin{erratas}
%     \errata{203}{3:23}{Se dice ..., pero debe decir ...}
%   \end{erratas}
% El primer argumento son los segundos (enlaza al minuto exacto del video),
% el segundo es la marca visible, el tercero es el texto.
\newenvironment{erratas}{\section*{Observaciones y erratas del video}\begin{itemize}}{\end{itemize}}
\newcommand{\errata}[3]{%
  \edef\errurl{https://www.youtube.com/watch?v=\videoid\string&t=#1}%
  \item[\textbf{\expandafter\href\expandafter{\errurl}{#2}}] #3}
) hace que las
% definiciones de la nota se llamen u.n, u.(n+1), ...; n es el primer número
% libre de esa unidad (ver la tabla de CONVENCIONES.md).
\newcommand{\unidad}[2]{\renewcommand{\thedefn}{#1.\arabic{defn}}\setcounter{defn}{\numexpr#2-1\relax}}

% Observaciones y erratas del video. Usar SOLO si hay algo que señalar:
%   \begin{erratas}
%     \errata{203}{3:23}{Se dice ..., pero debe decir ...}
%   \end{erratas}
% El primer argumento son los segundos (enlaza al minuto exacto del video),
% el segundo es la marca visible, el tercero es el texto.
\newenvironment{erratas}{\section*{Observaciones y erratas del video}\begin{itemize}}{\end{itemize}}
\newcommand{\errata}[3]{%
  \edef\errurl{https://www.youtube.com/watch?v=\videoid\string&t=#1}%
  \item[\textbf{\expandafter\href\expandafter{\errurl}{#2}}] #3}
.
% Compilar con LuaLaTeX:  lualatex nombre-de-la-nota.tex
\usepackage[spanish,provide=*]{babel}
\usepackage[T1]{fontenc}
\usepackage{lmodern}
\usepackage[a4paper,margin=2.6cm]{geometry}
\usepackage{amsmath,amssymb,amsthm}
\usepackage{xcolor}
\usepackage{tikz}
\usetikzlibrary{arrows.meta,babel}
\usepackage[colorlinks=true,urlcolor=blue!50!black,linkcolor=blue!50!black]{hyperref}

\hyphenpenalty=10000 \exhyphenpenalty=10000 \emergencystretch=3em

% Colores de las figuras (los mismos que usa la página)
\definecolor{extremo}{RGB}{47,107,79}
\definecolor{azul}{RGB}{44,74,154}

% Definiciones numeradas por unidad
\theoremstyle{definition}
\newtheorem{defn}{Definición}
% Numeración por unidad: \unidad{u}{n} (después de % Preámbulo compartido de las notas del curso de topología algebraica.
% Cada nota empieza con \documentclass[11pt]{article}, define \videoid{...}
% (el identificador del video de YouTube) y luego hace % Preámbulo compartido de las notas del curso de topología algebraica.
% Cada nota empieza con \documentclass[11pt]{article}, define \videoid{...}
% (el identificador del video de YouTube) y luego hace \input{preambulo}.
% Compilar con LuaLaTeX:  lualatex nombre-de-la-nota.tex
\usepackage[spanish,provide=*]{babel}
\usepackage[T1]{fontenc}
\usepackage{lmodern}
\usepackage[a4paper,margin=2.6cm]{geometry}
\usepackage{amsmath,amssymb,amsthm}
\usepackage{xcolor}
\usepackage{tikz}
\usetikzlibrary{arrows.meta,babel}
\usepackage[colorlinks=true,urlcolor=blue!50!black,linkcolor=blue!50!black]{hyperref}

\hyphenpenalty=10000 \exhyphenpenalty=10000 \emergencystretch=3em

% Colores de las figuras (los mismos que usa la página)
\definecolor{extremo}{RGB}{47,107,79}
\definecolor{azul}{RGB}{44,74,154}

% Definiciones numeradas por unidad
\theoremstyle{definition}
\newtheorem{defn}{Definición}
% Numeración por unidad: \unidad{u}{n} (después de \input{preambulo}) hace que las
% definiciones de la nota se llamen u.n, u.(n+1), ...; n es el primer número
% libre de esa unidad (ver la tabla de CONVENCIONES.md).
\newcommand{\unidad}[2]{\renewcommand{\thedefn}{#1.\arabic{defn}}\setcounter{defn}{\numexpr#2-1\relax}}

% Observaciones y erratas del video. Usar SOLO si hay algo que señalar:
%   \begin{erratas}
%     \errata{203}{3:23}{Se dice ..., pero debe decir ...}
%   \end{erratas}
% El primer argumento son los segundos (enlaza al minuto exacto del video),
% el segundo es la marca visible, el tercero es el texto.
\newenvironment{erratas}{\section*{Observaciones y erratas del video}\begin{itemize}}{\end{itemize}}
\newcommand{\errata}[3]{%
  \edef\errurl{https://www.youtube.com/watch?v=\videoid\string&t=#1}%
  \item[\textbf{\expandafter\href\expandafter{\errurl}{#2}}] #3}
.
% Compilar con LuaLaTeX:  lualatex nombre-de-la-nota.tex
\usepackage[spanish,provide=*]{babel}
\usepackage[T1]{fontenc}
\usepackage{lmodern}
\usepackage[a4paper,margin=2.6cm]{geometry}
\usepackage{amsmath,amssymb,amsthm}
\usepackage{xcolor}
\usepackage{tikz}
\usetikzlibrary{arrows.meta,babel}
\usepackage[colorlinks=true,urlcolor=blue!50!black,linkcolor=blue!50!black]{hyperref}

\hyphenpenalty=10000 \exhyphenpenalty=10000 \emergencystretch=3em

% Colores de las figuras (los mismos que usa la página)
\definecolor{extremo}{RGB}{47,107,79}
\definecolor{azul}{RGB}{44,74,154}

% Definiciones numeradas por unidad
\theoremstyle{definition}
\newtheorem{defn}{Definición}
% Numeración por unidad: \unidad{u}{n} (después de % Preámbulo compartido de las notas del curso de topología algebraica.
% Cada nota empieza con \documentclass[11pt]{article}, define \videoid{...}
% (el identificador del video de YouTube) y luego hace \input{preambulo}.
% Compilar con LuaLaTeX:  lualatex nombre-de-la-nota.tex
\usepackage[spanish,provide=*]{babel}
\usepackage[T1]{fontenc}
\usepackage{lmodern}
\usepackage[a4paper,margin=2.6cm]{geometry}
\usepackage{amsmath,amssymb,amsthm}
\usepackage{xcolor}
\usepackage{tikz}
\usetikzlibrary{arrows.meta,babel}
\usepackage[colorlinks=true,urlcolor=blue!50!black,linkcolor=blue!50!black]{hyperref}

\hyphenpenalty=10000 \exhyphenpenalty=10000 \emergencystretch=3em

% Colores de las figuras (los mismos que usa la página)
\definecolor{extremo}{RGB}{47,107,79}
\definecolor{azul}{RGB}{44,74,154}

% Definiciones numeradas por unidad
\theoremstyle{definition}
\newtheorem{defn}{Definición}
% Numeración por unidad: \unidad{u}{n} (después de \input{preambulo}) hace que las
% definiciones de la nota se llamen u.n, u.(n+1), ...; n es el primer número
% libre de esa unidad (ver la tabla de CONVENCIONES.md).
\newcommand{\unidad}[2]{\renewcommand{\thedefn}{#1.\arabic{defn}}\setcounter{defn}{\numexpr#2-1\relax}}

% Observaciones y erratas del video. Usar SOLO si hay algo que señalar:
%   \begin{erratas}
%     \errata{203}{3:23}{Se dice ..., pero debe decir ...}
%   \end{erratas}
% El primer argumento son los segundos (enlaza al minuto exacto del video),
% el segundo es la marca visible, el tercero es el texto.
\newenvironment{erratas}{\section*{Observaciones y erratas del video}\begin{itemize}}{\end{itemize}}
\newcommand{\errata}[3]{%
  \edef\errurl{https://www.youtube.com/watch?v=\videoid\string&t=#1}%
  \item[\textbf{\expandafter\href\expandafter{\errurl}{#2}}] #3}
) hace que las
% definiciones de la nota se llamen u.n, u.(n+1), ...; n es el primer número
% libre de esa unidad (ver la tabla de CONVENCIONES.md).
\newcommand{\unidad}[2]{\renewcommand{\thedefn}{#1.\arabic{defn}}\setcounter{defn}{\numexpr#2-1\relax}}

% Observaciones y erratas del video. Usar SOLO si hay algo que señalar:
%   \begin{erratas}
%     \errata{203}{3:23}{Se dice ..., pero debe decir ...}
%   \end{erratas}
% El primer argumento son los segundos (enlaza al minuto exacto del video),
% el segundo es la marca visible, el tercero es el texto.
\newenvironment{erratas}{\section*{Observaciones y erratas del video}\begin{itemize}}{\end{itemize}}
\newcommand{\errata}[3]{%
  \edef\errurl{https://www.youtube.com/watch?v=\videoid\string&t=#1}%
  \item[\textbf{\expandafter\href\expandafter{\errurl}{#2}}] #3}
) hace que las
% definiciones de la nota se llamen u.n, u.(n+1), ...; n es el primer número
% libre de esa unidad (ver la tabla de CONVENCIONES.md).
\newcommand{\unidad}[2]{\renewcommand{\thedefn}{#1.\arabic{defn}}\setcounter{defn}{\numexpr#2-1\relax}}

% Observaciones y erratas del video. Usar SOLO si hay algo que señalar:
%   \begin{erratas}
%     \errata{203}{3:23}{Se dice ..., pero debe decir ...}
%   \end{erratas}
% El primer argumento son los segundos (enlaza al minuto exacto del video),
% el segundo es la marca visible, el tercero es el texto.
\newenvironment{erratas}{\section*{Observaciones y erratas del video}\begin{itemize}}{\end{itemize}}
\newcommand{\errata}[3]{%
  \edef\errurl{https://www.youtube.com/watch?v=\videoid\string&t=#1}%
  \item[\textbf{\expandafter\href\expandafter{\errurl}{#2}}] #3}

\unidad{2}{6} % unidad 2; primer enunciado: 2.6

\title{Presentaciones de grupos}
\author{Luis Jorge Sánchez Saldaña\\[2pt]\small Facultad de Ciencias, UNAM}
\date{}

\begin{document}
\maketitle

\noindent\small Notas del video
\href{https://www.youtube.com/watch?v=V3OboR_tSEM}{\guillemotleft{}Presentaciones de grupos\guillemotright{}}
(curso de Introducción a la Topología Algebraica). Se ve que todo grupo es cociente de un grupo libre, se definen las presentaciones de grupos con varios ejemplos y se enuncia cómo definir homomorfismos desde un grupo dado por una presentación.
\normalsize

\section*{Todo grupo es cociente de un grupo libre}

Una \emph{presentación} es una forma de describir un grupo usando generadores y algunas reglas que hay que seguir para multiplicar los elementos del grupo, concretamente los generadores.

Sea $G$ un grupo. Recuerde que un subconjunto $S\subseteq G$ (no un subgrupo, solo un subconjunto) es un \emph{conjunto generador} de $G$ si $\langle S\rangle=G$. Esto se puede ver de dos formas: el subgrupo más pequeño que contiene a $S$ es todo $G$, o bien, todo elemento de $G$ se puede escribir como producto de elementos de $S$ y sus inversos. Para los fines de este video, la segunda descripción es la más útil.

Por otro lado, considerando a $S$ solo como conjunto, se puede formar el grupo libre $F_S$ generado por $S$, es decir, todas las palabras en el alfabeto formado por $S$ y sus inversos. Por la propiedad universal de los grupos libres, la inclusión $S\hookrightarrow G$ se extiende a un homomorfismo $\varphi\colon F_S\to G$, de modo que el siguiente diagrama conmuta:

\begin{center}
\begin{tikzpicture}[x=1cm,y=1cm]
  \node (S) at (0,1.6) {$S$};
  \node (G) at (2.2,1.6) {$G$};
  \node (F) at (0,0) {$F_S$};
  \draw[{Hooks[right,length=3pt]}-{Stealth[length=5pt]}] (S) -- (G);
  \draw[{Hooks[right,length=3pt]}-{Stealth[length=5pt]}] (S) -- (F);
  \draw[dashed,-{Stealth[length=5pt]}] (F) -- node[below right] {\small $\varphi$} (G);
\end{tikzpicture}
\end{center}

Además $\varphi$ es suprayectivo, porque su imagen es un subgrupo que contiene a $S$, y como $S$ es un conjunto generador, la imagen es todo $G$. Con esto se ha demostrado el siguiente teorema.

\begin{teo}\label{teo:cociente-libre}
Dado un grupo $G$, existe un grupo libre $F_S$ que se suprayecta en $G$. En particular, todo grupo es cociente de un grupo libre.
\end{teo}

\section*{Presentaciones}

Extrayendo la información del homomorfismo suprayectivo $\varphi\colon F_S\to G$, por el primer teorema de isomorfismo se tiene
\[
  G\cong F_S/\ker\varphi,
\]
y $\ker\varphi$ es un subgrupo normal de $F_S$. Si además se tiene un conjunto $R\subseteq\ker\varphi$ que genera normalmente al núcleo, es decir, tal que $\langle\!\langle R\rangle\!\rangle=\ker\varphi$, entonces
\[
  G\cong F_S/\ker\varphi=F_S/\langle\!\langle R\rangle\!\rangle.
\]

\begin{defn}[Presentación]\label{defn:presentacion}
En la situación anterior se dice que $G$ es el grupo generado por $S$ sujeto a las relaciones $R$, y se denota
\[
  G=\langle S\mid R\rangle.
\]
A esto se le llama una \emph{presentación} de $G$.
\end{defn}

Recapitulando, una presentación de $G$ es, esencialmente, un homomorfismo suprayectivo de un grupo libre en $G$ junto con un conjunto de generadores normales de su núcleo. Una presentación dice que todo elemento de $G$ se ve como producto de los elementos de $S$, es decir, como palabras en $S$, sujetas a las condiciones $R$. Como $R$ está contenido en $\ker\varphi$, que es un subgrupo de $F_S$, el conjunto $R$ también consiste de palabras en $S$.

\section*{Ejemplos}

\begin{enumerate}
  \item \emph{Grupos libres.} Para el grupo libre en un conjunto $S$ se tiene la presentación $F_S=\langle S\rangle$: la identidad $F_S\to F_S$ es suprayectiva y su núcleo es trivial, así que no hay ninguna relación. Cuando $R$ es vacío, simplemente no se escribe. Por ejemplo, si $F_2$ está generado por dos elementos $a,b$, entonces $F_2=\langle a,b\rangle$. Así, los grupos libres son grupos generados sin ninguna relación.

  Este ejemplo también muestra que las presentaciones no son únicas, porque hay muchas suprayecciones de grupos libres en un mismo grupo. Por ejemplo,
  \[
    F_2=\langle a,b\rangle=\langle a,b,c\mid c\rangle:
  \]
  se toma $F_3\to F_2$ con $a\mapsto a$, $b\mapsto b$ y $c\mapsto1$; el núcleo está generado (normalmente) por el elemento $c$, que se fue a $1$, y esa es la relación.

  \item \emph{El grupo abeliano libre en dos generadores.} Si los generadores de $\mathbb{Z}^2$ son $a$ y $b$ (que se pueden pensar como la base canónica $(1,0)$ y $(0,1)$), entonces
  \[
    \mathbb{Z}^2=\langle a,b\mid aba^{-1}b^{-1}\rangle.
  \]
  Otras formas muy comunes de escribirlo, aunque no son tal cual la notación, son $\langle a,b\mid ab=ba\rangle$ y $\langle a,b\mid aba^{-1}b^{-1}=1\rangle$, solo para especificar que ese elemento se va a $1$. Esto quiere decir que hay una suprayección $\varphi\colon F_2\to\mathbb{Z}^2$ cuyo núcleo está normalmente generado por el elemento $aba^{-1}b^{-1}$.

  \noindent\textbf{Ejercicio.} Verificar que el núcleo de $\varphi\colon F_2\to\mathbb{Z}^2$ está normalmente generado por $aba^{-1}b^{-1}$.

  \item \emph{El grupo abeliano libre en $n$ generadores.} De manera más general,
  \[
    \mathbb{Z}^n=\langle a_1,a_2,\dots,a_n\mid a_ia_j=a_ja_i\ \ \forall i,j\rangle,
  \]
  es decir, se impone que todos los generadores conmuten entre sí. Aunque se escribe una sola expresión, en realidad hay más de una relación, una por cada par $i,j$.

  \item \emph{El grupo cíclico de $n$ elementos.} El grupo $C_n$ está generado por un solo elemento, que es justamente la definición de cíclico. Entonces hay una suprayección $F_1=\mathbb{Z}\to C_n$ cuyo núcleo está generado por $n$ veces el generador, y así
  \[
    C_n=\langle a\mid a^n\rangle.
  \]

  \item \emph{Un ejemplo más interesante.} Sean $G_1=\langle S_1\mid R_1\rangle$ y $G_2=\langle S_2\mid R_2\rangle$. Nada impide definir el grupo
  \[
    G=\langle S_1\cup S_2\mid R_1\cup R_2\rangle;
  \]
  en realidad, cualquier conjunto de generadores y cualquier conjunto de relaciones dan un grupo. ?`Hay alguna forma de describir este grupo? Todo elemento de $G$ es una palabra en letras de $S_1$ y de $S_2$. Si hay varias letras de $S_1$ juntas, siempre se pueden aplicar las relaciones $R_1$, y lo mismo con las letras de $S_2$ y las relaciones $R_2$. Pero una letra de $S_1$ junto a una de $S_2$ no interactúan, porque las relaciones son palabras solo en $S_1$ o solo en $S_2$.

  \noindent\textbf{Ejercicio.} $\langle S_1\cup S_2\mid R_1\cup R_2\rangle\cong G_1*G_2$, el producto libre definido en un video anterior.
\end{enumerate}

\section*{Homomorfismos desde una presentación}

La siguiente proposición es consecuencia inmediata de la propiedad universal del cociente de grupos.

\begin{prop}\label{prop:hom-presentacion}
Sean $G=\langle S\mid R\rangle$, $K$ otro grupo y $f\colon S\to K$ una función (no un homomorfismo, una función cualquiera) tal que para toda relación $s_1\cdots s_m\in R$, con $s_i\in S$, se cumple
\[
  f(s_1)\cdots f(s_m)=1
\]
en $K$. Entonces existe un único homomorfismo $\bar f\colon G\to K$ tal que el siguiente diagrama conmuta:
\end{prop}

\begin{center}
\begin{tikzpicture}[x=1cm,y=1cm]
  \node (S) at (0,1.6) {$S$};
  \node (G) at (2.2,1.6) {$G$};
  \node (K) at (0,0) {$K$};
  \draw[-{Stealth[length=5pt]}] (S) -- (G);
  \draw[-{Stealth[length=5pt]}] (S) -- node[left] {\small $f$} (K);
  \draw[dashed,-{Stealth[length=5pt]}] (G) -- node[below right] {\small $\bar f$} (K);
\end{tikzpicture}
\end{center}

Es decir, para definir un homomorfismo de un grupo dado por una presentación con generadores $S$ y relaciones $R$, basta mandar los generadores y verificar que las relaciones van a dar al elemento $1$; por decirlo de alguna manera, que las relaciones se respeten. Esto sale de ver a $G$ como el cociente del grupo libre $F_S$ entre el subgrupo normalmente generado por $R$.

Esta es una introducción muy breve a las presentaciones de grupos, que se usarán para calcular grupos fundamentales con el teorema de van Kampen.

\begin{erratas}
  \errata{60}{1:00}{Se dice que todo elemento de $G$ es producto de elementos de $S$; debe decir producto de elementos de $S$ \emph{y de sus inversos} (así se escribe en estas notas).}
  \errata{523}{8:43}{Para que $\langle S_1\cup S_2\mid R_1\cup R_2\rangle\cong G_1*G_2$ hay que suponer que $S_1$ y $S_2$ son disjuntos (o tomar la unión disjunta); si comparten letras, el grupo que resulta no es el producto libre.}
  \errata{625}{10:25}{En la proposición se escriben las relaciones como $s_1\cdots s_m\in R$ con $s_i\in S$, pero las letras también pueden ser inversos de generadores (como en $aba^{-1}b^{-1}$): debe ser $s_i\in S\cup\bar S$, y si $s_i=s^{-1}$ con $s\in S$, entonces $f(s_i)$ se entiende como $f(s)^{-1}$.}
\end{erratas}

\end{document}
